\documentclass[11pt]{article}
\usepackage[a4paper,margin=25mm]{geometry}
\usepackage{amsmath,amssymb,amsthm,hyperref}
\newtheorem{theorem}{Theorem}
\newtheorem{lemma}{Lemma}
\newcommand{\PP}{\mathcal P}
\title{Rational box smoothing with polynomial-denominator conditional means}
\author{Proof independently reviewed by size\_bounds and independent\_directions}
\date{30 September 2026}
\begin{document}\maketitle
Throughout, ``polynomial'' means bounded by a fixed power of $m+1$,
with constants allowed to depend on fixed parameters but not on $m$.
The statement keeps the number of auxiliary rank features fixed.
It does not bound denominators of higher conditional moments.

\begin{theorem}\label{main}
Fix $h$. Suppose $K$ is polynomial in $m\ge2$, the nodes
$0<b_1<\cdots<b_K<1$ have gaps, including endpoint gaps, at least
an inverse polynomial, and $\psi_{\ell q}$, $1\le\ell\le h$, are
rational features with polynomially bounded absolute values and
polynomial common denominator. Suppose
\[
 T_{\ell j}=K^{-1}\sum_{q=1}^K\psi_{\ell q}b_q^j
 \quad(0\le j\le m)
\]
are rational, and that the $h$ numbers $T_{\ell1}$ have polynomial
common denominator. Finally, let $\phi_j$ be the orthonormal shifted
Legendre polynomials and assume the evaluation family
\[
 \psi_{\ell q}\phi_j(b_q),\qquad 1\le\ell\le h,\quad 0\le j<m,
\]
has empirical Gram matrix bounded below by an inverse polynomial.

Then there exist disjoint rational intervals
$I_q=[c_q-\epsilon,c_q+\epsilon]\subset(0,1)$, with positive
$\epsilon$ an inverse polynomial and with all their endpoints on
one polynomial rational grid, and rational numbers $a_q$ with
$|a_q|\le1/2$, such that the probability measure assigning each
$I_q$ mass $1/K$ and conditional density
\begin{equation}\label{density}
 \frac1{2\epsilon}\left[1+a_q\left(
       \frac{(x-c_q)^2}{\epsilon^2}-\frac13\right)\right]
\end{equation}
satisfies
\[
 \sum_q\psi_{\ell q}\int_{I_q}x^j\,d\beta(x)=T_{\ell j}
 \quad(0\le j\le m).
\]
Each box has conditional mean exactly $c_q$, hence these means
have a polynomial common denominator. All conditional moments are
rational. The centers can be made closer than any prescribed inverse
polynomial to the original nodes by adjusting the fixed exponents.
\end{theorem}
\begin{proof}
All spaces below have polynomial dimension. Polynomial evaluation,
first- and higher fixed-order derivative norms on $[0,1]$ are
polynomial in $m$; intervals will remain inside $[0,1]$.
We first replace the atoms by \emph{uniform} boxes and solve the
resulting moment problem over the reals, before rational rounding.

Choose rational $\epsilon=(m+1)^{-A}$ with $A$ a sufficiently large
fixed integer. For centers $c$, consider the nonconstant moment map
of the uniform boxes with feature weights $\psi_{\ell q}$.
Use the basis $f_j(t)=\int_0^t\phi_{j-1}(u)du$, $1\le j\le m$,
instead of powers. Its derivative at zero box width and centers $b$
is $K^{-1}$ times the stated evaluation family. The Gram hypothesis
gives a fixed right inverse with polynomial norm from Euclidean
residuals to maximum center displacement.
Uniform-box averaging changes each baseline moment by
$O(\operatorname{poly}(m)\epsilon^2)$, since its first centered
moment is zero. Its Jacobian differs by the same type of bound.
The usual fixed-inverse contraction therefore gives real centers
\[
 c_q^*=b_q+O(\operatorname{poly}(m)\epsilon^2)
\]
whose uniform boxes satisfy every target moment exactly. Choose $A$
large enough to preserve all node and endpoint gaps. The derivative
bounds in this contraction are polynomial; no denominator assertion
is made about $c^*$.

\medskip\noindent\emph{Rounding while preserving all first moments.}
The degree-one equations are exactly
\begin{equation}\label{linear}
 K^{-1}\sum_q\psi_{\ell q}c_q=T_{\ell1}.
\end{equation}
The feature matrix $(\psi_{\ell q})$ has rank $h$ and polynomial
Gram lower bound, as a subfamily of the assumed evaluation matrix.
Since $h$ is fixed, Cauchy--Binet supplies an $h\times h$ column
minor with polynomially bounded inverse. Round every $c_q^*$ to a
multiple of $1/L$, where $L=(m+1)^B$ with $B$ a sufficiently large
fixed integer, and correct only these $h$ selected centers to restore
\eqref{linear}. The resulting centers $c_q$ are rational and
\[
 \max_q|c_q-c_q^*|\le\operatorname{poly}(m)/L.
\]
They have a polynomial common denominator: the feature entries and
targets in~\eqref{linear} have polynomial common denominators, the
minor has fixed size, and Cramer's rule multiplies only a fixed number
of such denominators and bounded determinant numerators. Increasing
$B$ retains disjointness of the rational boxes and all endpoint gaps.

\medskip\noindent\emph{Correcting higher moments without moving the means.}
Use for the remaining equations the rationally rescalable basis
\[
 g_j(t)=\int_0^t(t-u)\phi_{j-2}(u)du,
       \qquad 2\le j\le m.
\]
Together with $1,t$, these span $\PP_m$. Add the density corrections
in~\eqref{density}; these preserve each box's mass and mean, because
for uniform $U$ on $[-1,1]$,
\[
 \mathbb E(U^2-1/3)=\mathbb E[U(U^2-1/3)]=0,
 \qquad \mathbb E[U^2(U^2-1/3)]=4/45.
\]
The response matrix $V$ of the remaining equations is exactly linear
in the $a_q$. Taylor expansion gives
\[
 V_{(\ell,j),q}
 =\frac{\psi_{\ell q}}K\mathbb E\bigl[
      g_j(c_q+\epsilon U)(U^2-1/3)\bigr]
 =\frac{2\epsilon^2}{45K}\psi_{\ell q}\phi_{j-2}(c_q)
       +O(\operatorname{poly}(m)\epsilon^4/K).
\]
The Gram hypothesis, center closeness, and sufficiently large $A$
therefore give a right inverse of $V$ with norm at most
$\epsilon^{-2}\operatorname{poly}(m)$. The moment residual after
center rounding is only $\operatorname{poly}(m)/L$, because the
uniform boxes at $c^*$ were exact. Consequently the linear system
has a solution with
\[
 \|a\|_\infty\le\operatorname{poly}(m)/(L\epsilon^2)\le1/2
\]
if $B$ is sufficiently large after $A$ is fixed.

The solution may be taken rational without sacrificing this bound.
Indeed, rescale each $g_j$ to remove its normalizing square root;
then $V$ and the residual vector are rational. The minimum Euclidean
norm solution $V^\top(VV^\top)^{-1}e$ is rational. It is unchanged
by any invertible rescaling of the equation rows, so the polynomial
bounds established in the orthonormal basis still apply.

The bracket in~\eqref{density} lies between $2/3$ and $4/3$, proving
positivity. It is even about $c_q$, so its conditional mean is $c_q$.
Rational coefficients and endpoints give rational conditional moments.
All required equations hold because the constant and linear ones were
preserved separately and the remaining spanning equations were solved.
Finally, increasing $A$ and then $B$ gives any fixed inverse-polynomial
center accuracy while all denominators claimed in the theorem remain
polynomial.
\end{proof}

\begin{theorem}[A fixed number of polynomial-denominator moments]\label{fixedmoments}
Fix integers $h,s\ge1$, and assume $m\ge s+1$. Under the hypotheses
of Theorem~\ref{main}, strengthen the arithmetic assumption by requiring
all $T_{\ell j}$, $1\le\ell\le h$, $1\le j\le s$, to have a
polynomial common denominator. Then there are rational boxes on one polynomial grid and positive
rational conditional densities of degree at most $s+1$, preserving all
the target weighted moments, for which \emph{all} conditional box moments
through degree $s$ have one polynomial common denominator. Each box has
mass $1/K$, and the center accuracy can be prescribed as before. The
conditional mean need not equal the box center in this extension. Density
bounds may be taken to be $1/2$ and $3/2$ relative to the uniform density
of the box.
\end{theorem}
\begin{proof}
First solve the real uniform-box problem exactly, obtaining the centers
$c_q^*$ as before. Let $\mu_{qj}^*$ be the conditional moment of degree
$j$ of the uniform box centered at $c_q^*$. Choose rational centers
$c_q$ on a fine common polynomial grid, and keep the common rational
halfwidth $\epsilon$. For each $1\le j\le s$, round all the numbers
$\mu_{qj}^*$ to a common polynomial grid and correct the same fixed
$h$ selected entries to impose
\[
 K^{-1}\sum_q\psi_{\ell q}\mu_{qj}=T_{\ell j},\qquad1\le\ell\le h.
\]
The fixed-size minor argument proves that all $\mu_{qj}$ have a
polynomial common denominator and differ from $\mu_{qj}^*$ by at most
$\operatorname{poly}(m)/L$. There are only the fixed number $hs$ of
additional corrected entries. Take $\mu_{q0}=1$.

In the uniform coordinate $u=(x-c_q)/\epsilon$, use the first $s$
Legendre modes to realize these conditional moments:
\[
 1+\sum_{l=1}^s a_{ql}P_l(u).
\]
For raw moment $j$, the response of mode $l$ is
$\mathbb E[(c_q+\epsilon U)^jP_l(U)]$. It vanishes when $j<l$,
and at $j=l$ equals $\epsilon^l\mathbb E[U^lP_l(U)]$, a nonzero
rational constant times $\epsilon^l$. This fixed-size triangular
system has rational inverse of norm at most $C_s\epsilon^{-s}$
for $0\le c_q\le1$. Thus it realizes the prescribed moments with
\[
 |a_{ql}|\le\operatorname{poly}(m)/(L\epsilon^s).
\]
They may be made as small as desired by increasing the polynomial
exponent in $L$.

It remains to restore all higher global moments. These low-mode
adjustments and center rounding have produced residual at most
$\operatorname{poly}(m)/(L\epsilon^s)$ in a polynomially bounded
basis. Add one further coefficient $a_{q,s+1}P_{s+1}(u)$ on each box.
It preserves every already fixed conditional moment of degree at
most $s$. For the remaining global equations use functions $g_j$,
$j=s+1,\ldots,m$, whose $(s+1)$st derivatives are respectively
$\phi_0,\ldots,\phi_{m-s-1}$ and whose first $s+1$ Taylor coefficients
at zero vanish. The response matrix divided by $\epsilon^{s+1}$
converges, with polynomial error $O(\operatorname{poly}(m)\epsilon)$,
to a nonzero rational constant times the evaluation subfamily
$K^{-1}\psi_{\ell q}\phi_{j-s-1}(b_q)$. Its right inverse therefore
has norm at most $\epsilon^{-(s+1)}\operatorname{poly}(m)$.
The rational minimum-norm solve, after harmless row rescaling, gives
\[
 |a_{q,s+1}|\le\operatorname{poly}(m)/(L\epsilon^{2s+1}).
\]
Choose first the fixed exponent in $\epsilon$ and then the one in $L$
large enough that $\sum_{l=1}^{s+1}|a_{ql}|\le1/2$ for every box.
Since $|P_l|\le1$ on $[-1,1]$, all densities have the asserted bounds.
Their endpoints and coefficients are rational, and their first $s$
conditional moments are the prescribed polynomial-denominator numbers.
Every required global moment is exact. All parameters are polynomial
for fixed $h,s$, as claimed.
\end{proof}

\paragraph{Arithmetic used by a synchronized family of applications.}
The fixed-minor rounding equations use $K T_{\ell j}$, rather than
$T_{\ell j}$ themselves. Thus if several applications use the same
rational feature minor, a common grid for the quantities $K T_{\ell j}$
and for all other rounded entries suffices to synchronize their output
low-moment denominators. This observation does not itself construct
such a shared minor or synchronize affine changes of interval.

\paragraph{Application to cumulative mass and first moment.}
The relevant fixed feature family is $\psi=(1,r_0,r_1)$, where $r_0$
is a rational cumulative mass and $r_1$ a rational cumulative first
moment. The trinary Legendre argument conditions this family when its
real baseline is
\[
 1,\quad R=t-\eta H_d,\quad M(t)=\int_0^t uR'(u)du.
\]
After sufficiently fine polynomial rational rounding, the empirical
Gram hypothesis remains valid. The targets of degree one are fixed
rational numbers in the intended uniform-moment construction. The
lemma thus produces rational pure boxes with polynomial-denominator
masses, endpoints, and means. Their higher moments may have much
larger denominators, which is deliberately outside the conclusion.
\end{document}
